\section{Introduction}
\label{sec:intro}
A distribution agreement in the Contract Understanding Atticus Dataset (CUAD) \citep{cuad} defines \emph{Trademarks} as only the names listed in Exhibit~A. It separately requires the distributor to report unauthorised use of the Trademarks. Asked whether a confusable but unlisted name must be reported, GPT-5.6-terra gives the wrong answer without the definition and the correct one with it (Table~\ref{tab:examples}). The notification clause alone does not expose the mistake: the ordinary meaning of \emph{trademark} suggests a broader duty. In our CUAD retrieval experiments, only 29\%--44\% of retrieved passage sets contain the needed definition. A system that checks an answer solely against those passages may therefore accept a fluent, source-inconsistent answer.

The example exposes a gap between \emph{using} evidence and \emph{knowing which evidence is missing}. Contracts, technical documentation, and communities all assign local scopes to terms \citep{lucy-bamman-glossaries}. A response may be consistent with the retrieved excerpt while conflicting with the full source. Checks for contradiction against the supplied passage will miss that particular omission; whether confidence can flag it is an empirical question. This matters for retrieval systems because improving a reader's willingness to follow context cannot recover a definition the retriever never supplied.

Knowledge-conflict studies usually provide competing evidence and ask which source the model follows \citep{longpre-etal-2021-entity,kc-survey,conflictbank,knowledge-conflict-faithfulness}. Here the relevant definition can be absent from the reader's input. On synthetic questions, supplying it repairs 209 errors and harms 6 previously correct answers. The immediate systems question is therefore which source-bound definitions deserve review and retrieval, not only whether a reader follows a definition it sees.

We construct diagnostic questions for term occurrences using the full source to propose answer keys, then score the keyed option with a small open model, with and without a specified definition. This moves source comparison to corpus preparation rather than asking a deployed reader to recognise an omission in its own input. The method's value depends on the questions and keys being sound, so we report a source review and test the limits of what the score predicts. Three findings carry the argument.

\begin{enumerate}
\item \textbf{Missing definitions can produce stable errors that confidence weakly detects} (Sections~\ref{sec:stable}--\ref{sec:detect}). On designed redefinitions, repeated sampling often preserves the wrong answer; the tested confidence signals are weakest on these terms. Evaluation should include omissions of local evidence, not only contradictions with supplied evidence.
\item \textbf{Source-linked scoring gives an offline review order} (Section~\ref{sec:predict}). A 7B proxy ranks stronger readers' errors on the same synthetic and contract questions, and concentrates estimated contract errors in a small part of a scored frame. It also tracks general item difficulty, so it is a triage score rather than a specific misreading detector or a proven future-request policy.
\item \textbf{The content and scope of a supplied definition determine its effect} (Section~\ref{sec:repair}). Source definitions strongly improve the constructed synthetic task, while a reader's prior gloss does little there. Drafted glosses can help or harm, and frequency remains a competitive default for deciding which definitions to curate first.
\end{enumerate}

